\section{Background and Related Work}\label{sec:related}
% Course rule: write this section yourselves. Organise by THEME and take a position; do not write one
% paragraph per paper. Cite only papers you have opened; take BibTeX from DBLP or the publisher.

\subsection{Blockchain academic credential systems}
\note{Themes, not lists. Which systems use a single issuing account? Which use permissioned networks or node
consensus? What do they assume about the honesty of the issuing employee? Candidates from your reading list:
EduCTX, CredenceLedger, Leka \& Selimi, DIAR (Rustemi et al.), Educert-chain (Rani et al.),
Cardenas-Quispe \& Pacheco, the IEEE Access systematic review, Al Hemairy et al.
Kazakhstan: Shakan et al., Ussatova et al., Begimbayeva et al.}

\subsection{Multi-party issuance in prior credential work}
\note{The closest baselines. Verify the \emph{exact} mechanism before describing it.
Lone \& Naaz (2020): is it really N-of-N multisignature?
Sarwar et al.\ (2026): node-level majority consensus in MultiChain, which differs from employee-level approval.
Tariq et al.\ (Cerberus): check its governance and issuance controls.
Nousias et al.\ (2022): issuance as an institutional workflow.
State plainly what these share and what none of them measures. That is your gap.}

\subsection{Role-based access control and separation of duty}
\note{RBAC foundations (Ferraiolo \& Kuhn 1992; Sandhu et al.\ 1996; the NIST model, 2000) and separation of
duty (Kuhn 1997). Our design uses static separation of duty (one address, one role) plus an approval-time
role-diversity constraint. Hyperledger Fabric endorsement policies already express role-aware AND/OR/threshold
rules. Acknowledge this and position the contribution as \emph{measurement} in the credential setting, not as
a new mechanism.}

\subsection{Credential standards}
\note{W3C Verifiable Credentials Data Model 2.0, Data Integrity 1.0 and Bitstring Status List: issuer, holder and
verifier roles, and revocation. Our commitment-plus-status design is a minimal on-chain analogue. Say what it
does not cover (selective disclosure, holder binding, DIDs).}
